\documentclass[instructions.tex]{subfiles}
\begin{document}
\chapter{Par le péché mortel on perd Dieu.}

\primo
Voici ce qu’il y a de plus funeste dans le péché mortel : il nous fait perdre Dieu et nous en sépare. La grâce sanctifiante nous unit à Dieu ; ainsi le péché mortel, en nous dépouillant de cette grâce, nous sépare de Dieu, et c’est dans cette séparation d’avec Dieu que consiste réellement la mort de l’âme ; parce que dans l’ordre surnaturel, Dieu est aussi véritablement la vie de notre âme, que l’âme dans l’ordre naturel est la vie du corps. Notre âme séparée de Dieu est donc aussi véritablement morte, dit saint Augustin, que le corps séparé de l’âme est véritablement mort.\footnote{Miserere animæ tuæ. \textit{Eccl.} 30. — Deus amissus est mors animæ, anima amissa est mors corporis. S. Aug.}

\secundo
Être séparé de Dieu, avoir perdu son Dieu, quelle perte ! Quand je perdrais tous les trésors de la terre, si je possède Dieu, rien n’est perdu pour moi ; mais quand je posséderais tous les empires de l’univers, si je perds Dieu, tout est perdu pour moi. Le bonheur des saints dans le ciel est d’être unis à Dieu et de le posséder. Le malheur des réprouvés dans l’enfer, c’est d’être séparés de Dieu et de l’avoir perdu.

Avoir perdu son Dieu, peut-il y avoir un plus grand malheur ? Oui, il y a un malheur plus grand : c’est après avoir perdu Dieu, de compter pour rien cette perte. Est-il possible que l’homme soit capable d’un pareil aveuglement ! Quoi, dit saint Augustin, vous pleurez sur un corps mort, parce que l’âme en est séparée ; et vous ne pleurez pas sur votre âme qui est séparée de Dieu.\footnote{Corpus luges à quo recessit anima ; animam non luges à qua recessit Deus. S. Aug.}

L’Écriture, parlant de Michas, à qui les voleurs avaient enlevé ses idoles, nous fait remarquer qu’il les redemandait avec des cris lamentables, et qu’étant interrogé quel était le sujet de ses larmes, il répondit : \textit{Hélas ! on m’a enlevé mes dieux, et vous me demandez pourquoi je pleure !} Si un idolâtre pleure amèrement la perte d’une idole, comment un chrétien peut-il être insensible sur la perte de Dieu ? On voit des personnes s’affliger pour les moindres disgrâces ; on entend dans des familles des emportements pour des pertes de rien, tandis qu’on s’y réjouit et qu’on est tranquille après avoir perdu Dieu par le crime. Ô péché, que tu pervertis étrangement l’esprit et le cœur de l’homme !

Saint Thomas d’Aquin, ce grand génie, ne pouvait comprendre comment on pouvait vivre après avoir perdu Dieu, et dormir dans sa disgrâce avec un péché mortel sur la conscience. Oh ! qu’on est misérable dans cet état ; et doublement misérable, si l’on n’en est pas touché ! \textit{Écoutez, ô mon peuple,} disait le Seigneur à Israël, \textit{voici ce qui montre la malice de votre cœur ; c’est qu’après m’avoir abandonné, vous avez encore étouffé tous les désirs de retour et tous les sentiments de crainte.} Jér. 2.

Si vous avez perdu Dieu, votre malheur est très-grand, il est incompréhensible ; mais il n’est pas sans remède. Allez faire aux pieds des ministres du Seigneur une humble confession de vos misères ; noyez vos péchés dans vos larmes, ou plutôt dans le sang de votre Rédempteur. Le repentir, l’amendement, les sacrements, sont les moyens qui peuvent vous rétablir et réparer la perte que vous avez faite.

\section{resolutions}

\primo Considérant le péché mortel comme donnant la mort à mon âme, je l’éviterai avec plus de soin que ce qui pourrait m’ôter la vie temporelle.

\secundo Et si j’ai jamais le malheur de me laisser entraîner dans un seul de ces misérables péchés, je me hâterai de sortir de cet état déplorable, en recourant promptement au sacrement de pénitence.

Ô mon Dieu, plutôt mourir mille fois que de me séparer encore de vous par un seul péché mortel !

\end{document}
